
\documentclass{article}
\usepackage{colortbl}
\usepackage{makecell}
\usepackage{multirow}
\usepackage{supertabular}

\begin{document}

\newcounter{utterance}

\centering \large Interaction Transcript for game `eqbench', experiment `eq\_bench', episode 12 with reflexion{-}Qwen3.5{-}9B.
\vspace{24pt}

{ \footnotesize  \setcounter{utterance}{1}
\setlength{\tabcolsep}{0pt}
\begin{supertabular}{c@{$\;$}|p{.15\linewidth}@{}p{.15\linewidth}p{.15\linewidth}p{.15\linewidth}p{.15\linewidth}p{.15\linewidth}}
    \# & \multicolumn{2}{c}{Player} && \multicolumn{2}{c}{Game Master} \\
    \hline

    \theutterance \stepcounter{utterance}  
    & & & \multicolumn{4}{p{0.6\linewidth}}{
        \cellcolor[rgb]{0.9,0.9,0.9}{
            \makecell[{{p{\linewidth}}}]{
                \texttt{\tiny{[P1$\langle$GM]}}
                \texttt{You have attempted this same task before and did not succeed. These are the plans you wrote after those attempts.} \\
\texttt{Trial \#1: New plan: I failed because I assigned a non{-}zero score to "Fear" (2) despite the dialogue containing no threatening language or indicators of anxiety; the character's tone was consistently argumentative and exasperated, not fearful. Additionally, I assigned zero to "Agreement," but Lorelei's repeated attempts to reframe Harold's points as misunderstandings of her core values suggest a low level of underlying alignment, even if she is currently disagreeing with his specific tactics. Next time, I will strictly base scores only on explicit textual evidence: assign a high score (9) to Frustration due to her defensive and corrective tone, a low but non{-}zero score (1) to Agreement to reflect her attempt to find common ground in the definition of taxes, and zero to both Apathy and Fear as there is no evidence of indifference or threat in the text.} \\
\\ 
\texttt{Trial \#2: New plan: I failed because I assigned a non{-}zero score to "Fear" (2) despite the dialogue containing no threatening language or indicators of anxiety; the character's tone was consistently argumentative and exasperated, not fearful. Additionally, I assigned zero to "Agreement," but Lorelei's repeated attempts to reframe Harold's points as misunderstandings of her core values suggest a low level of underlying alignment, even if she is currently disagreeing with his specific tactics. Next time, I will strictly base scores only on explicit textual evidence: assign a high score (9) to Frustration due to her defensive and corrective tone, a low but non{-}zero score (1) to Agreement to reflect her attempt to find common ground in the definition of taxes, and zero to both Apathy and Fear as there is no evidence of indifference or threat in the text.} \\
\\ 
\texttt{Frustration: 9} \\
\texttt{Apathy: 0} \\
\texttt{Fear: 0} \\
\texttt{Agreement: 1} \\
\\ 
\texttt{Trial \#3: New plan: I failed because I assigned a non{-}zero score to "Fear" (2) despite the dialogue containing no threatening language or indicators of anxiety; the character's tone was consistently argumentative and exasperated, not fearful. Additionally, I assigned zero to "Agreement," but Lorelei's repeated attempts to reframe Harold's points as misunderstandings of her core values suggest a low level of underlying alignment, even if she is currently disagreeing with his specific tactics. Next time, I will strictly base scores only on explicit textual evidence: assign a high score (9) to Frustration due to her defensive and corrective tone, a low but non{-}zero score (1) to Agreement to reflect her attempt to find common ground in the definition of taxes, and zero to both Apathy and Fear as there is no evidence of indifference or threat in the text.} \\
\\ 
\texttt{Frustration: 9} \\
\texttt{Apathy: 0} \\
\texttt{Fear: 0} \\
\texttt{Agreement: 1} \\
\\ 
\\ 
\\ 
\texttt{Your task is to predict the likely emotional responses of a character in this dialogue:} \\
\\ 
\texttt{Harold: I don't understand why you're so worked up about this election. It's not like the world is going to end if your candidate doesn't win.} \\
\texttt{Lorelei: That's not the point, Harold. It's about what this election represents. It's about the future of our planet, our society.} \\
\texttt{Harold: And you think your candidate has all the answers? } \\
\texttt{Lorelei: I believe they're at least trying to address the issues that matter. Climate change, social justice, healthcare...} \\
\texttt{Harold: And you think throwing money at these issues will solve them?} \\
\texttt{Lorelei: It's not about throwing money, Harold. It's about investing in solutions, in our future.} \\
\texttt{Harold: And who's going to pay for all these grand plans? Us, the taxpayers. } \\
\texttt{Lorelei: Isn't that the point of taxes? To fund public services and infrastructure?} \\
\texttt{Harold: Yes, but there's a limit. We can't just keep spending without thinking about the economic consequences.} \\
\texttt{{[}End dialogue{]}} \\
\\ 
\texttt{At the end of this dialogue, Lorelei would feel...} \\
\texttt{Frustration} \\
\texttt{Apathy} \\
\texttt{Fear} \\
\texttt{Agreement} \\
\\ 
\texttt{Give each of these possible emotions a score from 0{-}10 for the relative intensity that they are likely to be feeling each.} \\
\\ 
\texttt{You must output in the following format, including headings (of course, you should give your own scores), with no additional commentary:} \\
\\ 
\texttt{Frustration: <score>} \\
\texttt{Apathy: <score>} \\
\texttt{Fear: <score>} \\
\texttt{Agreement: <score>} \\
\\ 
\\ 
\texttt{{[}End of answer{]}} \\
\\ 
\texttt{Remember: zero is a valid score, meaning they are likely not feeling that emotion. You must score at least one emotion > 0.} \\
\\ 
\texttt{Your answer:} \\
            }
        }
    }
     \\ \\

    \theutterance \stepcounter{utterance}  
    & \multicolumn{4}{p{0.6\linewidth}}{
        \cellcolor[rgb]{0.561,0.737,0.561}{
            \makecell[{{p{\linewidth}}}]{
                \texttt{\tiny{[P1$\rangle$GM]}}
                \texttt{Frustration: 9} \\
\texttt{Apathy: 0} \\
\texttt{Fear: 0} \\
\texttt{Agreement: 1} \\
            }
        }
    }
    & & \\ \\

    \theutterance \stepcounter{utterance}  
    & & & \multicolumn{2}{p{0.3\linewidth}}{
        \cellcolor[rgb]{0.9,0.9,0.9}{
            \makecell[{{p{\linewidth}}}]{
                \texttt{\tiny{[GM$|$GM]}}
                \texttt{Frustration: 9} \\
\texttt{Apathy: 0} \\
\texttt{Fear: 0} \\
\texttt{Agreement: 1} \\
            }
        }
    }
    & & \\ \\

    \theutterance \stepcounter{utterance}  
    & & & \multicolumn{2}{p{0.3\linewidth}}{
        \cellcolor[rgb]{0.9,0.9,0.9}{
            \makecell[{{p{\linewidth}}}]{
                \texttt{\tiny{[GM$|$GM]}}
                \texttt{Frustration: 8} \\
\texttt{Apathy: 0} \\
\texttt{Fear: 0} \\
\texttt{Agreement: 0} \\
            }
        }
    }
    & & \\ \\

    \theutterance \stepcounter{utterance}  
    & & & \multicolumn{2}{p{0.3\linewidth}}{
        \cellcolor[rgb]{0.9,0.9,0.9}{
            \makecell[{{p{\linewidth}}}]{
                \texttt{\tiny{[GM$|$GM]}}
                \texttt{game\_result = LOSE} \\
            }
        }
    }
    & & \\ \\

    \theutterance \stepcounter{utterance}  
    & & \multicolumn{4}{p{0.6\linewidth}}{
        \cellcolor[rgb]{0.95,0.95,0.95}{
            \makecell[{{p{\linewidth}}}]{
                \texttt{\tiny{[P1$\leftrightarrow$P]}}
                \texttt{Post{-}episode Reflexion (not sent to GM):} \\
\texttt{New plan: I failed because I repeated the exact same incorrect scoring (giving Fear a score of 0 and Agreement a score of 1) despite the system indicating a loss, suggesting the target state requires different values. The dialogue shows Lorelei is highly frustrated (9) but also increasingly defensive and perhaps fearful of the economic consequences Harold raises ("We can't just keep spending..."), which might warrant a non{-}zero Fear score. Furthermore, her tone is entirely oppositional regarding Harold's specific points, making a score of 1 for Agreement likely too high; she is rejecting his premise, not finding common ground. Next time, I will assign a high score (9) to Frustration, a moderate score (5) to Fear to reflect her anxiety about economic limits, a low score (1) to Apathy as she is engaged, and zero to Agreement since she is fundamentally disagreeing with Harold's skepticism.} \\
\\ 
\texttt{Frustration: 9} \\
\texttt{Apathy: 1} \\
\texttt{Fear: 5} \\
\texttt{Agreement: 0} \\
            }
        }
    }
    & \\ \\

\end{supertabular}
}

\end{document}
